\subsection*{A short passage can hide a fact the reader already knows how to quote}
A source-pointing writer can return the actual shape of a historical glyph while composing a caption around it (Fig.~\ref{fig:pointing-causal-writer}). With the same pictured plate but opposite orders of the requested stages, it returned 188 of 192 captions exactly. Uniform pointing yielded zero and training with exchanged questions yielded one. Rephrasing the request on the same plates reduced accuracy to 21 of 192. The separate demonstration for the headword \textit{yán} (speech) reproduces an oracle form followed by a bronze form, then reverses them. The example illustrates which marks can appear in the answer. Its headword family was represented during training.

The same causal writer gave both exact facts for two of 128 simplified-Chinese CMRC passages and one of 128 traditional-Chinese DRCD passages. A transferred visual span reader gave 18 and 113, respectively. A learned choice between quoting and writing retained both sets of reading scores and all 188 familiar captions. Yet it selected writing for all 64 questions about named participants in short scenes. This gave zero exact names where quoting alone recovered 46. In the map scene, the first answer identified the mapmaker within a longer sentence. The second repeated the entire passage when asked who needed directions. Exact failure includes informative longer answers as well as selection of the wrong information.

Adding 1,738 short sentences from the existing natural training passages, each paired with its original question and answer, recovered the 46 names. Equal-length random excerpts recovered eleven. The sentence-trained selector returned 42 exact names on 64 further questions in sixteen new scenes, compared with four for random excerpts and zero for the original fit. Both name assignments and both roles were correct in six complete new scenes. The augmented selector quoted every answer in these two short-scene sets. The improvement came from selecting an available skill. The same change reduced exact answers on 571 open questions from eight to two, and none of 300 answers requiring new content ink was exact.

Poor document selection accounts for some missing evidence, but changing retrieval alone did little. Among 128 further open questions requiring a content glyph absent from the initial page, a learned ranker exposed the exact reference on five pages, compared with one for lexical retrieval. Even a selector allowed to see the reference exposed it only eleven times within the candidate windows. The writer matched none of the available human references under any of these selection rules. Exact-match failure alone cannot exclude a valid paraphrase. The short-scene results establish question-dependent selection of visible evidence, with transfer of linguistic knowledge from pretrained text models. They leave open how to generate a supported proposition beyond a quote or a familiar caption.
